\documentclass[10pt,a4paper]{amsart}
\usepackage{amsmath,amssymb,amsthm,graphicx}
% Retained minimal preamble for the Stacks Project chapter more-morphisms.tex (More on Morphisms).
% Verbatim excerpt of the chapter preamble: Xy-pic setup (native preamble lines 14-18) and the native
% theorem declarations with the mathematical macros (native preamble lines 154-217). Copied mechanically;
% nothing added, removed or redefined. The wrapper class and the amsmath/amssymb/amsthm/graphicx base
% packages are supplied by the assembler, so the class/externaldocument part of the native preamble is omitted.
% For commutative diagrams we use Xy-pic
\usepackage[all]{xy}

% We use 2cell for 2-commutative diagrams.
\xyoption{2cell}
% Theorem environments.
%
\theoremstyle{plain}
\newtheorem{theorem}[subsection]{Theorem}
\newtheorem{proposition}[subsection]{Proposition}
\newtheorem{lemma}[subsection]{Lemma}

\theoremstyle{definition}
\newtheorem{definition}[subsection]{Definition}
\newtheorem{example}[subsection]{Example}
\newtheorem{exercise}[subsection]{Exercise}
\newtheorem{situation}[subsection]{Situation}

\theoremstyle{remark}
\newtheorem{remark}[subsection]{Remark}
\newtheorem{remarks}[subsection]{Remarks}

\numberwithin{equation}{subsection}

% Macros
%
\def\lim{\mathop{\mathrm{lim}}\nolimits}
\def\colim{\mathop{\mathrm{colim}}\nolimits}
\def\Spec{\mathop{\mathrm{Spec}}}
\def\Hom{\mathop{\mathrm{Hom}}\nolimits}
\def\Ext{\mathop{\mathrm{Ext}}\nolimits}
\def\SheafHom{\mathop{\mathcal{H}\!\mathit{om}}\nolimits}
\def\SheafExt{\mathop{\mathcal{E}\!\mathit{xt}}\nolimits}
\def\Sch{\mathit{Sch}}
\def\Mor{\mathop{\mathrm{Mor}}\nolimits}
\def\Ob{\mathop{\mathrm{Ob}}\nolimits}
\def\Sh{\mathop{\mathit{Sh}}\nolimits}
\def\NL{\mathop{N\!L}\nolimits}
\def\CH{\mathop{\mathrm{CH}}\nolimits}
\def\proetale{{pro\text{-}\acute{e}tale}}
\def\etale{{\acute{e}tale}}
\def\QCoh{\mathit{QCoh}}
\def\Ker{\mathop{\mathrm{Ker}}}
\def\Im{\mathop{\mathrm{Im}}}
\def\Coker{\mathop{\mathrm{Coker}}}
\def\Coim{\mathop{\mathrm{Coim}}}

% Boxtimes
%
\DeclareMathSymbol{\boxtimes}{\mathbin}{AMSa}{"02}

%
% Macros for moduli stacks/spaces
%
\def\QCohstack{\mathcal{QC}\!\mathit{oh}}
\def\Cohstack{\mathcal{C}\!\mathit{oh}}
\def\Spacesstack{\mathcal{S}\!\mathit{paces}}
\def\Quotfunctor{\mathrm{Quot}}
\def\Hilbfunctor{\mathrm{Hilb}}
\def\Curvesstack{\mathcal{C}\!\mathit{urves}}
\def\Polarizedstack{\mathcal{P}\!\mathit{olarized}}
\def\Complexesstack{\mathcal{C}\!\mathit{omplexes}}
% \Pic is the operator that assigns to X its picard group, usage \Pic(X)
% \Picardstack_{X/B} denotes the Picard stack of X over B
% \Picardfunctor_{X/B} denotes the Picard functor of X over B
\def\Pic{\mathop{\mathrm{Pic}}\nolimits}
\def\Picardstack{\mathcal{P}\!\mathit{ic}}
\def\Picardfunctor{\mathrm{Pic}}
\def\Deformationcategory{\mathcal{D}\!\mathit{ef}}
\newenvironment{reference}{}{}
\newenvironment{slogan}{}{}
\newenvironment{history}{}{}
\usepackage{cleveref}
\crefname{originallemma}{original lemma}{original lemmas}
\Crefname{originallemma}{Original lemma}{Original lemmas}
\makeatletter
\newlabel{algebra-lemma-henselization-different}{{algebra-lemma-henselization-different}{source}{}{sourcepointer-1}{}}
\newlabel{algebra-lemma-henselization-different@cref}{{[originallemma][0][]algebra-lemma-henselization-different}{[1][0][]source}{}{}{}}
\makeatother
\crefname{originallemma}{original lemma}{original lemmas}
\Crefname{originallemma}{Original lemma}{Original lemmas}
\makeatletter
\newlabel{algebra-lemma-local-flat-ff}{{algebra-lemma-local-flat-ff}{source}{}{sourcepointer-2}{}}
\newlabel{algebra-lemma-local-flat-ff@cref}{{[originallemma][0][]algebra-lemma-local-flat-ff}{[1][0][]source}{}{}{}}
\makeatother
\crefname{originallemma}{original lemma}{original lemmas}
\Crefname{originallemma}{Original lemma}{Original lemmas}
\makeatletter
\newlabel{algebra-lemma-strict-henselization-different}{{algebra-lemma-strict-henselization-different}{source}{}{sourcepointer-3}{}}
\newlabel{algebra-lemma-strict-henselization-different@cref}{{[originallemma][0][]algebra-lemma-strict-henselization-different}{[1][0][]source}{}{}{}}
\makeatother
\crefname{originallemma}{original lemma}{original lemmas}
\Crefname{originallemma}{Original lemma}{Original lemmas}
\makeatletter
\newlabel{lemma-nr-minimal-primes}{{lemma-nr-minimal-primes}{source}{}{sourcepointer-4}{}}
\newlabel{lemma-nr-minimal-primes@cref}{{[originallemma][0][]lemma-nr-minimal-primes}{[1][0][]source}{}{}{}}
\makeatother
\crefname{originallemma}{original lemma}{original lemmas}
\Crefname{originallemma}{Original lemma}{Original lemmas}
\makeatletter
\newlabel{morphisms-lemma-etale-at-point}{{morphisms-lemma-etale-at-point}{source}{}{sourcepointer-5}{}}
\newlabel{morphisms-lemma-etale-at-point@cref}{{[originallemma][0][]morphisms-lemma-etale-at-point}{[1][0][]source}{}{}{}}
\makeatother
\crefname{originallemma}{original lemma}{original lemmas}
\Crefname{originallemma}{Original lemma}{Original lemmas}
\makeatletter
\newlabel{properties-definition-number-of-branches}{{properties-definition-number-of-branches}{source}{}{sourcepointer-6}{}}
\newlabel{properties-definition-number-of-branches@cref}{{[originallemma][0][]properties-definition-number-of-branches}{[1][0][]source}{}{}{}}
\makeatother
\crefname{originallemma}{original lemma}{original lemmas}
\Crefname{originallemma}{Original lemma}{Original lemmas}
\makeatletter
\newlabel{properties-lemma-number-of-branches-1}{{properties-lemma-number-of-branches-1}{source}{}{sourcepointer-7}{}}
\newlabel{properties-lemma-number-of-branches-1@cref}{{[originallemma][0][]properties-lemma-number-of-branches-1}{[1][0][]source}{}{}{}}
\makeatother
\title{Controlling the number of branches of the fibres by the normalisation}
\author{The Stacks Project Authors}
\date{stacks-project:more-morphisms.tex:local-snapshot-20261011; GNU FDL 1.2 or later}
\begin{document}
\maketitle
\noindent\textit{Editorial note.} Statement span is lines 10333-10352 of sources/stacks-more-morphisms/pinned-original-source.tex; proof span is lines 10354-10402, the complete author proof (47 body lines).\par\medskip
\noindent\textit{Editorial note.} Navigation stubs: 7 cited-but-not-included labels are supplied as non-load-bearing pointers with their pinned tags where the snapshot tag index has them.\par\medskip
\noindent\textit{Editorial note.} extra\_packages is empty on purpose: the retained preamble already loads xy with options and a second declaration clashes in the assembler wrapper.\par\medskip
\noindent\textit{Editorial note.} Text is retained verbatim, including the author's own typos and placeholders.\par\medskip
\noindent\textit{Allowed standard prerequisites (editorial).}\begin{itemize}
\item Explicit prerequisite boundary. Normalisation of a locally Noetherian scheme, finiteness of the normalisation, and the behaviour of fibres of a finite morphism; the spreading argument is retained.
\end{itemize}\medskip
\ifdefined\hypertarget\hypertarget{sourcepointer-1}{}\fi
\noindent\textit{Source reference.} Original lemma algebra-lemma-henselization-different: sources/stacks-more-morphisms/pinned-original-source.tex, source lines 10402--10402. Cited by the retained statement or proof but not part of the collected spans; supplied only as a navigation pointer. The argument does not depend on its content.\par\medskip
\ifdefined\hypertarget\hypertarget{sourcepointer-2}{}\fi
\noindent\textit{Source reference.} Original lemma algebra-lemma-local-flat-ff: sources/stacks-more-morphisms/pinned-original-source.tex, source lines 10402--10402. Cited by the retained statement or proof but not part of the collected spans; supplied only as a navigation pointer. The argument does not depend on its content.\par\medskip
\ifdefined\hypertarget\hypertarget{sourcepointer-3}{}\fi
\noindent\textit{Source reference.} Original lemma algebra-lemma-strict-henselization-different: sources/stacks-more-morphisms/pinned-original-source.tex, source lines 10402--10402. Cited by the retained statement or proof but not part of the collected spans; supplied only as a navigation pointer. The argument does not depend on its content.\par\medskip
\ifdefined\hypertarget\hypertarget{sourcepointer-4}{}\fi
\noindent\textit{Source reference.} Original lemma lemma-nr-minimal-primes: sources/stacks-more-morphisms/pinned-original-source.tex, source lines 10303--10309. Cited by the retained statement or proof but not part of the collected spans; supplied only as a navigation pointer. The argument does not depend on its content.\par\medskip
\ifdefined\hypertarget\hypertarget{sourcepointer-5}{}\fi
\noindent\textit{Source reference.} Original lemma morphisms-lemma-etale-at-point: sources/stacks-more-morphisms/pinned-original-source.tex, source lines 10402--10402. Cited by the retained statement or proof but not part of the collected spans; supplied only as a navigation pointer. The argument does not depend on its content.\par\medskip
\ifdefined\hypertarget\hypertarget{sourcepointer-6}{}\fi
\noindent\textit{Source reference.} Original lemma properties-definition-number-of-branches: sources/stacks-more-morphisms/pinned-original-source.tex, source lines 10402--10402. Cited by the retained statement or proof but not part of the collected spans; supplied only as a navigation pointer. The argument does not depend on its content.\par\medskip
\ifdefined\hypertarget\hypertarget{sourcepointer-7}{}\fi
\noindent\textit{Source reference.} Original lemma properties-lemma-number-of-branches-1: sources/stacks-more-morphisms/pinned-original-source.tex, source lines 10402--10402. Cited by the retained statement or proof but not part of the collected spans; supplied only as a navigation pointer. The argument does not depend on its content.\par\medskip
\begin{lemma}
\label{lemma-nr-branches}
Let $X$ be a scheme and $x \in X$ a point. Then
\begin{enumerate}
\item the number of branches of $X$ at $x$ is equal to
the supremum of the number of irreducible components of $U$
passing through $u$ taken over elementary \'etale neighbourhoods
$(U, u) \to (X, x)$,
\item the number of geometric branches of $X$ at $x$ is equal to
the supremum of the number of irreducible components of $U$
passing through $u$ taken over \'etale neighbourhoods
$(U, u) \to (X, x)$,
\item $X$ is unibranch at $x$ if and only if for every
elementary \'etale neighbourhood $(U, u) \to (X, x)$ there
is exactly one irreducible component of $U$ passing through $u$, and
\item $X$ is geometrically unibranch at $x$ if and only if for every
\'etale neighbourhood $(U, u) \to (X, x)$ there
is exactly one irreducible component of $U$ passing through $u$.
\end{enumerate}
\end{lemma}

\begin{proof}
Parts (3) and (4) follow from parts (1) and (2) via
Properties, Lemma \ref{properties-lemma-number-of-branches-1}.

\medskip\noindent
Proof of (1). Let $\Spec(A)$ be an affine open neighbourhood
of $x$ and let $\mathfrak p \subset A$ be the prime ideal
corresponding to $x$. We may replace $X$ by $\Spec(A)$ and
it suffices to consider affine elementary \'etale neighbourhoods
$(U, u)$ in the supremum as they form a cofinal subsystem.
Recall that the henselization $A_\mathfrak p^h$
is the colimit of the rings $B_\mathfrak q$ over the category
of pairs $(B, \mathfrak q)$ where $B$ is an \'etale $A$-algebra
and $\mathfrak q$ is a prime lying over $\mathfrak p$ with
$\kappa(\mathfrak q) = \kappa(\mathfrak p)$, see
Algebra, Lemma \ref{algebra-lemma-henselization-different}.
These pairs $(B, \mathfrak q)$ correspond exactly to
the affine elementary \'etale neighbourhoods $(U, u)$
by the correspondence between rings and affine schemes.
Observe that irreducible components of $\Spec(B)$
passing through $\mathfrak q$ are exactly the minimal
prime ideals of $B_\mathfrak q$. The number of minimal
primes of $A_\mathfrak p^h$ is the number of branches
of $X$ at $x$ by Properties, Definition
\ref{properties-definition-number-of-branches}.
Observe that the transition
maps $B_\mathfrak q \to B'_{\mathfrak q'}$ in the system
are all flat. Since a flat local ring map is faithfully flat
(Algebra, Lemma \ref{algebra-lemma-local-flat-ff})
we see that the lemma follows
from Lemma \ref{lemma-nr-minimal-primes}.

\medskip\noindent
Proof of (2). The proof is the same as the proof of (1), except that we use
Algebra, Lemma \ref{algebra-lemma-strict-henselization-different}.
There is a tiny difference: given a separable algebraic closure
$\kappa^{sep}$ of $\kappa(x)$ for every \'etale neighbourhood
$(U, u)$ we can choose a $\kappa(x)$-embedding
$\phi : \kappa(u) \to \kappa^{sep}$
because $\kappa(u)/\kappa(x)$ is finite separable
(Morphisms, Lemma \ref{morphisms-lemma-etale-at-point}).
Hence we can look at the supremum over all triples
$(U, u, \phi)$ where $(U, u) \to (X, x)$ is an affine
\'etale neighbourhood and $\phi : \kappa(u) \to \kappa^{sep}$
is a $\kappa(x)$-embedding. These triples correspond
exactly to the triples in
Algebra, Lemma \ref{algebra-lemma-strict-henselization-different}
and the rest of the proof is exactly the same.
\end{proof}

\end{document}
